\section{Restricted fixed-rank sums in a prime cyclic group}
\label{app:restricted-fold}

We record the additive-combinatorial input used by the local response theorem.  The proof is the standard polynomial-method argument, included to keep the arithmetic dependency explicit.

\begin{theorem}[Restricted fixed-rank fold]\label{thm:restricted-fold}
Let $S$ be a cyclic group of prime order $p$, let $A\subseteq S$ with $|A|=r$, and let $0\le h\le r$.  If
\[
\Sigma_h:\operatorname{FinSub}_h(A)\longrightarrow S,
\qquad O\longmapsto\sum_{a\in O}a,
\]
then
\[
|\operatorname{im}\Sigma_h|
\ge \min\{p,h(r-h)+1\}.
\]
In particular, if $h(r-h)+1\ge p$, then $\Sigma_h$ is surjective.
\end{theorem}

\begin{proof}
The endpoint cases $h=0,r$ are immediate, so assume $0<h<r$.  Choose a proof-local identification $S\simeq\mathbf F_p$; the final statement is invariant under this choice.  Since $A\subseteq\mathbf F_p$, one has $r\le p$.

Put
\[
M=\min\{p-1,h(r-h)\}.
\]
Suppose for contradiction that $|\operatorname{im}\Sigma_h|\le M$.  Enlarge the image to a set $C\subseteq\mathbf F_p$ of cardinality exactly $M$.  Consider
\[
P(x_1,\ldots,x_h)
=
\prod_{c\in C}\left(\sum_{i=1}^h x_i-c\right)
\prod_{1\le i<j\le h}(x_j-x_i).
\tag{RF1}\label{eq:rf-polynomial}
\]
The polynomial vanishes on $A^h$: repeated coordinates kill the Vandermonde factor, while distinct coordinates have sum in the image of $\Sigma_h$, hence in $C$.

We use two elementary coefficient facts.  First, if a polynomial $F\in K[x_1,\ldots,x_h]$ has total degree at most $\sum_i d_i$, if
\[
[x_1^{d_1}\cdots x_h^{d_h}]F\ne0,
\]
and if finite sets $A_i$ satisfy $|A_i|>d_i$, then $F$ is nonzero somewhere on $\prod_iA_i$.  This follows by tensoring the one-variable Lagrange coefficient-extraction functionals.

Second, for distinct nonnegative integers $k_1,\ldots,k_h$ and
\[
m=\sum_i k_i-\frac{h(h-1)}2,
\]
one has
\[
[x_1^{k_1}\cdots x_h^{k_h}]
(x_1+\cdots+x_h)^m
\prod_{i<j}(x_j-x_i)
=
m!\frac{\prod_{i<j}(k_j-k_i)}{\prod_i k_i!}.
\tag{RF2}\label{eq:rf-vandermonde}
\]
Indeed, expand the Vandermonde determinant; the remaining alternating sum is a determinant of falling factorials, hence the ordinary Vandermonde determinant.

For every $0\le m\le h(r-h)$ one may choose distinct
\[
0\le k_1<\cdots<k_h\le r-1
\]
with
\[
\sum_i k_i=m+\frac{h(h-1)}2.
\tag{RF3}\label{eq:rf-exponents}
\]
To see this, write $R=r-h$ and $m=qh+s$ with $0\le s<h$.  Take a nondecreasing sequence $\lambda_i\in[0,R]$ with $h-s$ copies of $q$ and $s$ copies of $q+1$, and put $k_i=i-1+\lambda_i$.  If $q=R$, then $m\le hR$ forces $s=0$, so the displayed bounds remain valid.

It remains to connect the coefficient in \eqref{eq:rf-vandermonde} with the actual polynomial \eqref{eq:rf-polynomial}.  Put
\[
H_C(T)=\prod_{c\in C}(T-c),
\qquad
V(x)=\prod_{i<j}(x_j-x_i).
\]
Because $|C|=M$, $H_C$ is monic of degree $M$, so
\[
H_C(T)=T^M+R_C(T),
\qquad \deg R_C<M.
\tag{RF4}\label{eq:rf-top-homogeneous}
\]
The Vandermonde $V$ is homogeneous of degree $h(h-1)/2$.  Hence every monomial of
\[
R_C(x_1+\cdots+x_h)V(x)
\]
has total degree strictly below
\[
M+\frac{h(h-1)}2.
\]
Consequently, for every multi-index $k$ of exactly that total degree,
\[
[x^k]P
=
[x^k](x_1+\cdots+x_h)^M V(x).
\tag{RF5}\label{eq:rf-coefficient-bridge}
\]
Thus no lower homogeneous term of $H_C$ can contribute to the coefficient used below.

Apply \eqref{eq:rf-exponents} with $m=M$.  The corresponding monomial has total degree $M+h(h-1)/2$, so by \eqref{eq:rf-coefficient-bridge} and \eqref{eq:rf-vandermonde} its coefficient in $P$ is
\[
M!\frac{\prod_{i<j}(k_j-k_i)}{\prod_i k_i!}.
\tag{RF6}\label{eq:rf-nonzero-coefficient}
\]
This is nonzero in $\mathbf F_p$: one has $M\le p-1$ and $0\le k_i\le r-1\le p-1$, so all factorials in numerator and denominator are nonzero; moreover the distinct integers $k_i\in\{0,\ldots,p-1\}$ have pairwise differences nonzero modulo $p$.  The coefficient detector therefore contradicts the vanishing of $P$ on $A^h$.

Hence $|\operatorname{im}\Sigma_h|>M$, which is exactly
\[
|\operatorname{im}\Sigma_h|\ge\min\{p,h(r-h)+1\}.
\]
Transporting the conclusion back across the proof-local identification $S\simeq\mathbf F_p$ shows that no basis or generator survives in the theorem statement.
\end{proof}
